\chapter{Metodo de Diferenças Finitas}
\label{cap:MDF}

\section{Contextualização do Método}

O método de diferenças finitas é o método mais antigo para soluções de equações diferenciais parciais (PDEs) e normalmente utilizado para geometrias simples. O ponto de partida do método se baseia na própria definição de derivada:
\begin{equation}
    f'(x) = \lim_{h \to 0} \frac{f(x+h) - f(x)}{h}
\end{equation}
Dado qualquer x em que esse limite exista, f' é a derivada de f, sendo o valor de f' em x a inclinação da reta tangente à f no ponto (x,f(x)). De acordo com Ferziger Peric (2002), o método de diferenças finitas (MDF) para aproximação de derivadas parciais, possui três maneiras de aproximação: diferença progressiva, diferença regressiva e a terceira diferença central. A tabela \ref{methodAppr} apresenta a diferença entre os três métodos.
\begin{table}[h!]
\label{methodAppr}
\centering
\begin{tabular}{|c|c|}
\hline
\textbf{Método} & \textbf{Aproximação} \\ \hline

Forward Difference 
& 
$\displaystyle f'(x) \approx \frac{f(x+h) - f(x)}{h}$ 
\\ \hline

Backward Difference 
& 
$\displaystyle f'(x) \approx \frac{f(x) - f(x-h)}{h}$ 
\\ \hline

Central Difference 
& 
$\displaystyle f'(x) \approx \frac{f(x+h) - f(x-h)}{2h}$ 
\\ \hline

\end{tabular}
\end{table}

Dentre essas, a utilização da diferença central é a que melhor aproxima a solução real. Na figura~\ref{devappr}, Ferziger e Perić (2002) apresentam as três aproximações para a derivada de primeira ordem; observa-se que a diferença central é a que melhor descreve a tangente de $f(x)$.

\begin{figure}
\label{devappr}
    \centering
    \includegraphics[width=0.5\linewidth]{image.png}
    \caption {Enter Caption}
\end{figure}

\section{Séries de Taylor}

\subsection{Definição em uma variável}

Seja $f: I \subset \mathbb{R} \to \mathbb{R}$ uma função infinitamente 
derivável em um intervalo $I$ contendo o ponto $x_0$. A série de Taylor de 
$f$ em torno de $x_0$ é formalmente dada por
\begin{equation}
    f(x) \sim \sum_{n=0}^{\infty} \frac{f^{(n)}(x_0)}{n!}(x - x_0)^n,
\end{equation}
onde $f^{(n)}$ denota a $n$-ésima derivada de $f$. Dizemos que a série 
\emph{converge} para $f$ em um ponto $x$ se a soma infinita converge e o 
valor da soma é precisamente $f(x)$.

Na prática numérica, truncamos a série após um certo número de termos, 
obtendo o \emph{polinômio de Taylor} de ordem $p$:
\begin{equation}
    T_p(x; x_0) = \sum_{n=0}^{p} \frac{f^{(n)}(x_0)}{n!}(x - x_0)^n.
\end{equation}
Escrevemos então
\begin{equation}
    f(x) = T_p(x; x_0) + R_{p+1}(x),
\end{equation}
onde $R_{p+1}(x)$ é o \emph{resto de Taylor}, que representa o erro de 
aproximação ao truncar a série no termo de ordem $p$.

\subsection{Forma do resto e ordem de aproximação}

Uma forma clássica do resto (forma de Lagrange) é
\begin{equation}
    R_{p+1}(x) = \frac{f^{(p+1)}(\xi)}{(p+1)!}(x - x_0)^{p+1},
\end{equation}
para algum $\xi$ entre $x$ e $x_0$. Se tomamos $x = x_0 + h$, com $h$ 
pequeno, obtemos
\begin{equation}
    R_{p+1}(x_0 + h) = \mathcal{O}(h^{p+1}),
\end{equation}
isto é, o erro de truncamento é proporcional a $h^{p+1}$ quando $h \to 0$. 

Na notação de \emph{ordem de grandeza}, escrever
\begin{equation}
    f(x_0 + h) = T_p(x_0 + h; x_0) + \mathcal{O}(h^{p+1})
\end{equation}
significa que existe uma constante $C$ tal que
\begin{equation}
    |R_{p+1}(x_0 + h)| \le C |h|^{p+1}
\end{equation}
para $h$ suficientemente pequeno. Essa é a base para dizer que 
um esquema numérico é de ordem $p$: o erro de truncamento decresce como 
uma potência $p$ do parâmetro de malha $h$.

Em problemas bidimensionais, costuma-se usar $h$ como um parâmetro genérico 
de malha, por exemplo $h = \max\{h_x,h_y\}$, onde $h_x$ e $h_y$ são os passos 
em cada direção.

\subsection{Séries de Taylor como ferramenta conceitual}

Nesse contexto, a série de Taylor é sobretudo uma ferramenta para descrever 
o comportamento local da função. Ao aproximarmos $f$ por um polinômio, 
substituímos uma função possivelmente complicada por um objeto simples de 
manipular analiticamente e numericamente. Essa descrição local permite 
aproximar derivadas por combinações lineares de valores de $f$, analisar a 
precisão dessas aproximações e deduzir a ordem de um método numérico, 
compreendendo como o erro se comporta quando refinamos a malha. Em outras 
palavras, a expansão de Taylor fornece uma espécie de ``zoom'' local sobre 
a função, a partir do qual se organizam, de maneira sistemática, as fórmulas 
de diferenças finitas e a estimativa de seus erros.

\section{Diferenças finitas a partir de séries de Taylor}

\subsection{Discretização unidimensional}

Para fixar ideias, considere uma função $f(x)$ definida em $[a,b]$. 
Discretizamos o intervalo em pontos igualmente espaçados:
\begin{equation}
    x_i = a + i h, \quad i = 0,1,2,\dots,N,
\end{equation}
onde $h = (b-a)/N$ é o passo da malha. Em vez da função contínua $f(x)$, 
passamos a trabalhar com os valores discretos $f_i = f(x_i)$.

As fórmulas clássicas de diferenças finitas para derivadas vêm diretamente 
das expansões de Taylor. A seguir, apresentamos as aproximações em 1D que 
serão reutilizadas na construção de esquemas em malhas 2D.

\subsection{Diferença progressiva para a primeira derivada}

Considere os pontos $x_i$ e $x_{i+1} = x_i + h$. Expandindo $f$ em torno 
de $x_i$:
\begin{equation}
    f(x_i + h) = f(x_i) + h f'(x_i) + \frac{h^2}{2} f''(x_i)
                 + \frac{h^3}{6} f^{(3)}(x_i) + \cdots.
\end{equation}
Subtraindo $f(x_i)$ em ambos os lados e dividindo por $h$:
\begin{equation}
    \frac{f(x_i + h) - f(x_i)}{h}
    = f'(x_i) + \frac{h}{2} f''(x_i) + \frac{h^2}{6} f^{(3)}(x_i) + \cdots.
\end{equation}
Se ignoramos os termos de ordem superior em $h$, obtemos a aproximação
\begin{equation}
    f'(x_i) \approx \frac{f(x_{i+1}) - f(x_i)}{h},
\end{equation}
chamada de \emph{diferença progressiva}. O desenvolvimento acima mostra que
\begin{equation}
    f'(x_i) = \frac{f(x_{i+1}) - f(x_i)}{h} + \mathcal{O}(h),
\end{equation}
isto é, a fórmula é de \emph{primeira ordem}.

\subsection{Diferença regressiva e central}

De forma análoga, obtemos a \emph{diferença regressiva}
\begin{equation}
    f'(x_i) \approx \frac{f(x_i) - f(x_{i-1})}{h}
    \quad \text{com erro } \mathcal{O}(h),
\end{equation}
e a \emph{diferença central}
\begin{equation}
    f'(x_i) \approx \frac{f(x_{i+1}) - f(x_{i-1})}{2h}
    \quad \text{com erro } \mathcal{O}(h^2).
\end{equation}

\subsection{Diferença central para a segunda derivada em 1D}

Para a segunda derivada, a combinação das expansões em $x_i \pm h$
leva à fórmula central clássica
\begin{equation}
    f''(x_i) \approx \frac{f(x_{i+1}) - 2f(x_i) + f(x_{i-1})}{h^2},
\end{equation}
com erro de truncamento local $\mathcal{O}(h^2)$. Essa expressão será o 
bloco básico para a discretização de derivadas segundas em cada direção de 
uma malha bidimensional.

\subsection{Extensão para uma malha bidimensional}

Agora consideramos uma função $u(x,y)$ definida em um retângulo
\begin{equation}
    \Omega = (a,b) \times (c,d).
\end{equation}
Introduzimos uma malha cartesiana uniforme:
\begin{equation}
    x_i = a + i h_x, \quad i = 0,1,\dots,N_x,
    \qquad
    y_j = c + j h_y, \quad j = 0,1,\dots,N_y,
\end{equation}
onde $h_x = (b-a)/N_x$ e $h_y = (d-c)/N_y$. A aproximação discreta de 
$u(x,y)$ é dada pelos valores
\begin{equation}
    u_{i,j} \approx u(x_i,y_j).
\end{equation}

A partir das fórmulas unidimensionais, obtemos aproximações para as 
derivadas parciais de segunda ordem em cada direção. Em $x$:
\begin{equation}
    u_{xx}(x_i,y_j) \approx 
    \frac{u_{i+1,j} - 2u_{i,j} + u_{i-1,j}}{h_x^2},
\end{equation}
com erro de truncamento $\mathcal{O}(h_x^2)$, e, em $y$:
\begin{equation}
    u_{yy}(x_i,y_j) \approx 
    \frac{u_{i,j+1} - 2u_{i,j} + u_{i,j-1}}{h_y^2},
\end{equation}
com erro de truncamento $\mathcal{O}(h_y^2)$.

Para operadores que envolvem a soma dessas derivadas (como o Laplaciano),
a combinação dessas expressões leva naturalmente ao \emph{esquema em malha 2D}
que será utilizado no exemplo de equação elíptica.

\section{Erro de truncamento, consistência, estabilidade e convergência}

\subsection{Erro local de truncamento}

Seja $\mathcal{L}$ um operador diferencial atuando sobre uma função suave 
$u$ e $\mathcal{L}_h$ um operador discreto, dependente do parâmetro de malha 
$h$ (em 2D, pode-se tomar $h = \max\{h_x,h_y\}$). Em um ponto de malha 
$(x_i,y_j)$, o \emph{erro local de truncamento} é definido por
\begin{equation}
    \tau_h(x_i,y_j) 
    = \mathcal{L}[u](x_i,y_j) - \mathcal{L}_h[u](x_i,y_j).
\end{equation}
Usando expansões de Taylor em torno de $(x_i,y_j)$, escreve-se
\begin{equation}
    \tau_h(x_i,y_j) = C\,h^p + \mathcal{O}(h^{p+1}),
\end{equation}
para algum $p \ge 1$ e constante $C$ que depende de derivadas de ordem mais 
alta de $u$. Nesse caso, dizemos que o esquema é de \emph{ordem $p$}. De 
forma mais compacta, costuma-se escrever
\begin{equation}
    \tau_h(x_i,y_j) = \mathcal{O}(h^p).
\end{equation}

A ordem do erro de truncamento fornece a taxa com que o erro numérico tende 
a zero quando a malha é refinada. Em problemas bem postos, essa informação 
se traduz na taxa de convergência da solução discreta, desde que o esquema 
seja estável.

\subsection{Consistência, estabilidade e convergência}

Considere um problema contínuo, por exemplo um problema de valor de contorno,
\begin{equation}
    \mathcal{L}[u](x,y) = g(x,y), \quad (x,y) \in \Omega,
\end{equation}
com condições de contorno apropriadas em $\partial\Omega$. A discretização 
leva a um problema em malha:
\begin{equation}
    \mathcal{L}_h[u_h](x_i,y_j) = g(x_i,y_j),
\end{equation}
para os nós internos $(x_i,y_j)$, acompanhado de condições de contorno 
discretizadas nos nós da borda. Denotamos por $u_h$ a solução numérica 
obtida pelo método de diferenças finitas em uma malha de passo $h$ 
(ou $(h_x,h_y)$ em duas dimensões), e por $u$ a solução exata do 
problema contínuo.

Três conceitos são centrais nessa análise. O primeiro é o de \textbf{consistência}. 
Dizemos que o esquema é consistente se o erro de truncamento $\tau_h(x_i,y_j)$ 
tende a zero quando $h \to 0$, isto é, se $\mathcal{L}_h[u](x_i,y_j)$ se aproxima 
de $\mathcal{L}[u](x_i,y_j)$ para funções suaves $u$. Em termos de Taylor, isso 
é garantido quando o erro de truncamento é de ordem positiva em $h$ (por exemplo, 
$\mathcal{O}(h)$, $\mathcal{O}(h^2)$ e assim por diante).

O segundo conceito é o de \textbf{estabilidade}. Um esquema é estável quando 
perturbações nos dados (incluindo erros de arredondamento) não são amplificadas 
de forma descontrolada pelo processo numérico. Em problemas elípticos 
discretizados em malhas bidimensionais, a estabilidade costuma estar associada 
ao fato de a matriz do sistema linear ser bem condicionada e satisfazer 
propriedades como um princípio do máximo discreto. Já em problemas evolutivos 
(do tipo $u_t = \mathcal{L}[u]$), a estabilidade em esquemas explícitos 
frequentemente impõe restrições do tipo CFL (\emph{Courant--Friedrichs--Lewy}), 
que ligam o passo de tempo $\Delta t$ aos passos espaciais $h_x$ e $h_y$, 
surgindo condições típicas do tipo
\[
   \Delta t \le C \,\min\{h_x^2,h_y^2\},
\]
para alguma constante $C>0$ dependente do problema.

Por fim, temos a \textbf{convergência}. Um esquema é convergente se a solução 
numérica $u_h$ se aproxima da solução exata $u$ quando $h \to 0$ (e, em problemas 
dependentes do tempo, também quando $\Delta t \to 0$). Em problemas lineares bem 
postos, um resultado fundamental garante que, sob hipóteses adequadas, 
consistência e estabilidade em conjunto implicam convergência: em termos 
informais, se o esquema reproduz corretamente a equação e não amplifica 
descontroladamente os erros, então a solução discreta tende à solução contínua.

As expansões de Taylor são a ferramenta padrão para verificar a consistência 
e determinar a ordem de um esquema em malhas 2D. Já a estabilidade é tratada 
por técnicas específicas ao tipo de equação (elíptica, parabólica, hiperbólica) 
e ao formato da discretização (explícita, implícita ou mista).

\section{Exemplo: equação elíptica bidimensional}

\subsection{Formulação contínua e condições de contorno}

Considere o problema de valor de contorno bidimensional
\begin{equation}
    -\left(u_{xx}(x,y) + u_{yy}(x,y)\right) = f(x,y), 
    \quad (x,y) \in \Omega = (a,b) \times (c,d),
\end{equation}
com condições de contorno de Dirichlet prescritas na borda:
\begin{equation}
    u(x,y) = g(x,y), \quad (x,y) \in \partial\Omega.
\end{equation}
O objetivo é aproximar numericamente $u(x,y)$ em uma malha 2D.

\subsection{Discretização do domínio 2D}

Discretizamos $\Omega$ em uma malha cartesiana uniforme:
\begin{equation}
    x_i = a + i h_x, \quad i = 0,1,\dots,N_x,
    \qquad
    y_j = c + j h_y, \quad j = 0,1,\dots,N_y,
\end{equation}
e definimos
\begin{equation}
    u_{i,j} \approx u(x_i,y_j), 
    \qquad
    f_{i,j} = f(x_i,y_j), 
    \qquad
    g_{i,j} = g(x_i,y_j).
\end{equation}

Os nós de contorno são aqueles para os quais $i=0$ ou $i=N_x$, ou ainda 
$j=0$ ou $j=N_y$. Neles, as condições de Dirichlet impõem diretamente:
\begin{equation}
    u_{i,j} = g_{i,j}, \quad \text{para } (i,j) \text{ na borda da malha}.
\end{equation}

Os nós internos são aqueles com $1 \le i \le N_x-1$ e 
$1 \le j \le N_y-1$, onde a equação diferencial deve ser aproximada.

\subsection{Aproximação do Laplaciano em 2D}

No interior da malha, aproximamos as derivadas segundas por diferenças centrais:
\begin{align}
    u_{xx}(x_i,y_j) &\approx 
    \frac{u_{i+1,j} - 2u_{i,j} + u_{i-1,j}}{h_x^2}, \\
    u_{yy}(x_i,y_j) &\approx 
    \frac{u_{i,j+1} - 2u_{i,j} + u_{i,j-1}}{h_y^2}.
\end{align}
Substituindo na equação contínua, obtemos, para cada nó interno $(i,j)$:
\begin{equation}
    -\left(
      \frac{u_{i+1,j} - 2u_{i,j} + u_{i-1,j}}{h_x^2}
      + 
      \frac{u_{i,j+1} - 2u_{i,j} + u_{i,j-1}}{h_y^2}
    \right)
    = f_{i,j}.
\end{equation}

Reorganizando os termos em $u_{i,j}$:
\begin{equation}
    \left(
      \frac{2}{h_x^2} + \frac{2}{h_y^2}
    \right) u_{i,j}
    - \frac{1}{h_x^2}\left(u_{i+1,j} + u_{i-1,j}\right)
    - \frac{1}{h_y^2}\left(u_{i,j+1} + u_{i,j-1}\right)
    = f_{i,j}.
\end{equation}

Quando um dos vizinhos $(i\pm1,j)$ ou $(i,j\pm1)$ está na borda da malha, 
o valor correspondente de $u$ é conhecido pela condição de contorno e pode 
ser movido para o lado direito da equação. O resultado é um sistema linear 
para os valores desconhecidos $u_{i,j}$ nos nós internos.

Pelas expansões de Taylor em cada direção, temos
\begin{equation}
    u_{xx}(x_i,y_j) 
      = \frac{u_{i+1,j} - 2u_{i,j} + u_{i-1,j}}{h_x^2}
        + \mathcal{O}(h_x^2),
    \qquad
    u_{yy}(x_i,y_j) 
      = \frac{u_{i,j+1} - 2u_{i,j} + u_{i,j-1}}{h_y^2}
        + \mathcal{O}(h_y^2),
\end{equation}
de modo que o erro de truncamento local do esquema para o Laplaciano é
\begin{equation}
    \tau_h(x_i,y_j) = \mathcal{O}(h_x^2 + h_y^2).
\end{equation}
Se $h_x$ e $h_y$ são da mesma ordem de grandeza, podemos escrever 
simplesmente $\tau_h(x_i,y_j) = \mathcal{O}(h^2)$, o que mostra que o 
esquema é de segunda ordem em malhas 2D.

\subsection{Comentário sobre estabilidade e convergência}

Para esse problema elíptico estacionário com condições de Dirichlet, a 
discretização por diferenças centrais em malha cartesiana produz uma matriz 
simétrica definida positiva, associada a um operador discreto que preserva 
um princípio do máximo em sentido apropriado. Essas propriedades estão 
relacionadas à estabilidade do esquema: pequenos erros nos dados ou nas 
condições de contorno não produzem oscilações artificiais de grande 
amplitude na solução discreta.

Como o esquema é consistente de ordem dois e estável, e o problema contínuo 
é bem posto, obtém-se convergência da solução numérica $u_h$ para a solução 
exata $u$, com erro global tipicamente da ordem de $h^2$ em normas adequadas.


\section{Condições de contorno em malhas bidimensionais}

Para obter uma solução única para o problema contínuo e para a sua formulação 
discreta, é necessário especificar condições de contorno no bordo do domínio. 
Em uma malha 2D, isso significa prescrever informações nos nós que pertencem
às linhas e colunas de fronteira, por exemplo $i=0$, $i=N_x$, $j=0$ e $j=N_y$.

As duas classes mais comuns de condições de contorno são:

\begin{itemize}
    \item \emph{Condições de Dirichlet}: o valor da variável é conhecido 
    no contorno. Em termos discretos, temos
    \[
    u_{i,j} = u^{\text{bc}}_{i,j} 
    \quad \text{para } (x_i,y_j) \in \partial\Omega,
    \]
    onde $u^{\text{bc}}_{i,j}$ é o valor prescrito no contorno.
    
    \item \emph{Condições de Neumann}: o gradiente normal da variável é 
    conhecido no contorno. Em uma fronteira $x=a$, por exemplo, temos
    \[
    \pd{u}{n} = \pd{u}{x} = g(y) \quad \text{em } x=a,
    \]
    e, na forma discreta, pode-se utilizar uma fórmula de diferença finita 
    apropriada, como
    \[
    \left(\pd{u}{x}\right)_{0,j}
    \approx \frac{u_{1,j} - u_{0,j}}{h_x} = g(y_j),
    \]
    ou esquemas de ordem mais alta, dependendo do problema.
\end{itemize}

Quando a condição de contorno envolve derivadas e não há ponto de malha 
imediatamente fora do domínio, é comum introduzir pontos fictícios 
(\emph{ghost points}) para fechar o sistema de equações. Nesses casos, 
a condição de contorno é discretizada de modo a relacionar o valor no 
ponto fictício com os valores internos, eliminando explicitamente a 
variável extra.

\subsection{Condições de contorno no interior do sistema discreto}

Segundo Ferziger e Perić (2002), uma vez conhecidos os valores da variável 
na fronteira (Dirichlet) ou o seu gradiente (Neumann), essas informações 
são incorporadas diretamente nas equações de diferenças finitas dos pontos 
de fronteira ou de seus vizinhos imediatos. Na prática, isso significa:

\begin{itemize}
    \item substituir os valores prescritos em Dirichlet diretamente nas 
    equações (ou no vetor de termos conhecidos);
    \item no caso de Neumann, reescrever a derivada normal por uma 
    aproximação por diferenças finitas e, a partir dela, obter uma relação
    discreta que entra como equação adicional ou que elimina um valor
    de contorno em favor de valores internos.
\end{itemize}



\section{Sistema de equações algébricas}

A aproximação por diferenças finitas produz, em geral, uma equação algébrica
para cada ponto interno da malha. Para problemas lineares, o resultado da 
discretização pode ser escrito na forma
\begin{equation}
\label{system}
A_p \,u_p + \sum_{l} A_l \,u_l = Q_p,
\end{equation}
onde:
\begin{itemize}
    \item $p$ representa o ponto central da \emph{molécula computacional}
    (por exemplo $(i,j)$ em 2D);
    \item o índice $l$ percorre os pontos vizinhos que aparecem na 
    aproximação (por exemplo, norte, sul, leste, oeste na fórmula de cinco pontos);
    \item $A_p$ e $A_l$ são coeficientes que dependem dos parâmetros físicos 
    (como coeficientes de difusão e convecção) e dos passos de malha;
    \item $Q_p$ reúne termos fonte e contribuições das condições de contorno.
\end{itemize}

No caso de uma molécula computacional de cinco pontos para o Laplaciano em 
uma malha uniforme com $h_x = h_y = h$, a equação em um ponto interno $(i,j)$ 
pode ser escrita como
\begin{equation}
\label{system_laplace_2d}
- u_{i-1,j} - u_{i+1,j} - u_{i,j-1} - u_{i,j+1} + 4u_{i,j}
= h^2 f_{i,j},
\end{equation}
onde $f_{i,j} \approx f(x_i,y_j)$ é o termo fonte avaliado no ponto 
$(x_i,y_j)$. Comparando com \eqref{system}, vemos que $A_p = 4$, os 
coeficientes dos vizinhos são $A_l = -1$ e $Q_p = h^2 f_{i,j}$ (ajustado por 
eventuais contribuições das condições de contorno).

Reunindo as equações de todos os pontos internos, obtemos um sistema linear
esparso, que pode ser escrito em forma matricial como
\begin{equation}
\label{system_matrix}
A\,\bm{u} = \bm{Q},
\end{equation}
onde $\bm{u}$ é o vetor que contém os valores desconhecidos 
$u_{i,j}$ em todos os pontos internos. O número de incógnitas coincide
com o número de equações geradas pela discretização, resultando em um 
sistema determinado.



\section{Síntese conceitual}

A conexão entre séries de Taylor e o método de diferenças finitas pode ser 
resumida em alguns pontos centrais. Em primeiro lugar, há a ideia de 
\textbf{localidade}: séries de Taylor descrevem o comportamento local da 
função em torno de um ponto, e diferenças finitas aproveitam essa descrição 
para aproximar derivadas a partir de valores em nós vizinhos. Em seguida, 
vem a forma como essas fórmulas são construídas: as expressões de diferenças 
finitas nada mais são do que combinações lineares de valores da função 
cujo objetivo é cancelar, nas expansões de Taylor, termos de baixa ordem em 
$h$, isolando a derivada desejada e empurrando o erro para potências mais 
altas de $h$.

Outro aspecto importante é a \textbf{ordem do método}. A potência dominante 
de $h$ no erro de truncamento define a ordem do esquema. Em malhas 2D, essa 
ordem leva em conta o comportamento conjunto dos passos $h_x$ e $h_y$, mas, 
quando ambos têm a mesma escala, é comum falar em uma única ordem global em 
termos de um parâmetro $h$ típico da malha.

Do ponto de vista analítico, a tríade \textbf{consistência, estabilidade e 
convergência} organiza a discussão: a análise via Taylor fornece a 
consistência e a ordem de truncamento, enquanto a estabilidade depende da 
natureza da equação (elíptica, parabólica, hiperbólica) e do tipo de 
esquema (explícito, implícito, misto). Em muitos problemas lineares bem 
postos, saber que o esquema é consistente e estável é suficiente para 
garantir que a solução numérica converge para a solução exata quando a malha 
é refinada.

Por fim, vale destacar a \textbf{generalização para malhas multidimensionais}. 
As ideias discutidas no caso unidimensional se estendem de forma natural a 
malhas 2D e 3D. No exemplo apresentado, a combinação das fórmulas de segunda 
ordem em cada direção levou ao esquema clássico de cinco pontos para o 
Laplaciano em duas dimensões, amplamente utilizado na solução numérica de 
problemas elípticos. Séries de Taylor fornecem, assim, a linguagem e a 
ferramenta analítica que permitem construir, interpretar e avaliar esquemas 
de diferenças finitas em malhas bidimensionais para aproximação de derivadas 
e solução numérica de equações diferenciais parciais.